\scanpage{213}
% archive sheet number: 212

\mspage{109}

cette opération \unsure{ci|de $Z$} soit transitive -- soit donc $\varphi_0\in\Pi_Z$, \unsure{déjà|disons}, et $\varphi\in\Pi_Z$, on peut écrire

\[
\varphi=g_0^\vee\varphi_0=\varphi_0g_0^{-1}\qquad\text{où }(g_0,g_0^\vee)\in\operatorname{Aut}(\Gamma),
\]

\inserted{\uline{est-il dans $Z_\Gamma$ ?}} uniquement déterminé, on \deleted{doit} exprimer la

\[
\text{condition}\qquad\varphi\in\Pi_Z\quad\text{i.e.}
\]

% slanted
\margin{\uline{Lemme} Soit $\varphi_0\in\Pi_Z$ une $Z$-jauge, et $g\in\operatorname{Aut}(\Gamma)$, telle que l'on ait \deleted{\unsure{$g$}} $\underbrace{g_S^\vee\varphi_0=\varphi_0g_S^{\unsure{-1|{}}}}_{\varphi}$ (i.e. $(g_S,g_S^\vee)$ définit une homothétie de la structure de graphe $A_{\varphi_0}$ associée à $\varphi_0$). Conditions équivalentes : a) $g_S^\vee\varphi'=\varphi'g_S^{-1}$ pour \uline{tt} $\varphi'$ de la forme $T_\lambda\varphi_0$ ($\lambda\in Z$) i.e. $g_S$ est une homothétie \inserted{pour \uline{ttes} les str. de graphes $A_\varphi$} \inserted{$\leftarrow$ \uline{tt} $\varphi'\in\Pi$} b) $\varphi$ ($\overset{\text{déf}}=g_S^\vee\varphi_0=\varphi_0g_S^{-1}$) est également une $Z$-jauge c) $g\in\operatorname{Cent}(Z)$}

\[
\underbrace{\lambda^\vee\varphi\lambda}=\underbrace{\varphi},\quad\forall(\lambda,\lambda^\vee)\in Z\quad\text{i.e., \unsure{écrivant|écrivons} }\varphi=\varphi_0g_0^{-1}
\]

\[
\deleted{\unsure{\lambda^\vee g_0^\vee\varphi_0\lambda=g_0^\vee\varphi_0}}\quad\unsure{2}
\]

\[
\begin{aligned}
&\deleted{\unsure{g_0}}\quad\lambda^\vee\varphi_0g_0^{-1}\lambda=\varphi_0g_0^{-1}\qquad\text{i.e.}\\
&\underbrace{\lambda^\vee\varphi_0}_{\varphi_0\lambda^{-1}}g_0^{-1}(\lambda)=\varphi_0\qquad\text{i.e.}\\
&\deleted{\unsure{R}}\qquad\deleted{\varphi_0=\lambda_0}\\
&\lambda^{-1}g_0^{-1}(\lambda)=1\\
&\lambda=g_0^{-1}(\lambda)\qquad\text{i.e. } g_0\ \deleted{\text{normalise}}\ \inserted{\text{centralise}}\ \text{tout }\lambda\in\unsure{\mathfrak S|\mathscr S}_S\\
&\text{ou }g_0(\lambda)=\lambda\qquad\text{provenant de }\deleted{\unsure{\mathfrak S_S}}\ Z
\end{aligned}
\]

\deleted{donc} on trouve

\[
\left\{\begin{aligned}
&\text{l'ens. }\Pi_Z\ \text{des jauges de }\Gamma=(S,S^\vee,R)\\
&\text{compatibles avec le groupe d'autom.}\\
&Z\subset\operatorname{Aut}(\Gamma),\ \text{i.e satisfaisant }\lambda^\vee\varphi\lambda\quad\forall(\lambda,\lambda^\vee)\in Z,\\
&\text{est un torseur sous le \uline{centralisateur}}\\
&\text{de }Z\ \text{dans }\operatorname{Aut}(\Gamma),\ \text{opérant par}\\
&T_g\varphi_0=g_{S^\vee}\varphi_0=\varphi_0g_S^{-1}
\end{aligned}\right. \tag*{(68)}
\]

\margin{\uline{faux}}

% slanted
\margin{Il faut écrire la petite vérification que si on a $g_S^\vee\varphi_{\unsure{a|0}}=\varphi g_S^{-1}$ pour \uline{un} $\varphi\in\Pi$, on l'a pour tous (\unsure{ceci pour|dès que} $g\in\unsure{\mathcal G|G}$).}

Ouf -- il aura fallu de la sueur, pour accoucher de cette petite souris !

\unsure{tout ce|tout-ce|tendre} qu'on a montré c'est que l'application

\[
\deleted{\unsure{\varphi}}\ g\longmapsto g_{S^\vee}\varphi_0=\varphi_0g_S^{-1}
\]

est une \uline{bijection} du \deleted{le} sous-ensemble $\mathcal G(\varphi_0)$ de $\mathcal G$, formé des $g$ satisfaisant $g_{S^\vee}\varphi_0=\varphi_0g_S^{-1}$ i.e. qui sont des $A_{\varphi_0}$-homothéties (\uline{et qui ne forment pas \inserted{néc.} un groupe} !) avec l'ens $\Pi_Z$ des $Z$-jauges !

% slanted
\margin{\uline{NB} on a $\mathcal G(\varphi_0)=\{g\in\operatorname{Aut}(\Gamma)\mid\varphi_0(g_S)g^\vee=1\}$ i.e. $g^\vee=\varphi_0(g_S^{-1})$}

% slanted
\margin{Il est donc prudent d'étendre un voile pudique sur l'ens $\Pi_Z$, qui n'est pas bien compris et peut-être \unsure{déshérité|d'intérêt|dénué d'intérêt}, pour se borner décidément, aux ens $\Pi$}

\uline{Mais il est vrai que $Z$ opère librement sur $\Pi_Z$}
